\subsection{INNER PRODUCTS: CASE OF ALGEBRAS}

Let \(E = \bigoplus_{p \geqslant 0} E_p\) be a graded A-algebra of type N and P a graded A-module of type Z; for every homogeneous element \(x \in E_p\) , left multiplication by \(x\) is an A-linear mapping \(e(x)\) of E into itself which is graded of degree \(p\) . For every element \(u \in \text{Homgr}_A(E, P)\) , the right inner product of \(u\) by \(x\) , denoted by \(u \llcorner x\) , is the element \(u \circ e(x)\) of \(\text{Homgr}_A(E, P)\) . We also write \((i(x))(u) = u \llcorner x\) and we see that \(i(x)\) is a graded endomorphism of degree \(p\) of the graded A-module \(\text{Homgr}_A(E, P)\) . If now \(x = \sum_{p \geqslant 0} x_p\) is an arbitrary element of E (with \(x_p \in E_p\) for all \(p \geqslant 0\) , \(x_p = 0\) except for a finite number of values of \(p\) ), we write \(i(x) = \sum_{p=0}^{\infty} i(x_p)\) , which is therefore an endomorphism of the A-module \(\text{Homgr}_A(E, P)\) .

To remember which element, in the expression \(u \llcorner x\) , ``operates'' on the other, observe that the element x which ``operates'' on u is placed at the free end of the horizontal line in \(\llcorner\).

The associativity of the algebra E goes over to the relation \(e(xy) = e(x) \circ e(y)\) for x, y homogeneous; whence, by definition of \(i(x)\) ,

\[
i (x y) = i (y) \circ i (x)\tag{35}
\]

first for \(x, y\) homogeneous and then, by linearity, for arbitrary \(x, y\) in E; this may also be written

\[
(u \llcorner x) \llcorner y = u \llcorner (x y)\tag{36}
\]

for \(x, y\) in \(\mathbf{E}\) and \(u \in \operatorname{Homgr}_{\mathbf{A}}(\mathbf{E}, \mathbf{P})\) ; as on the other hand clearly \(i(1)\) is the identity mapping (since this follows from \(e(1) = 1_{\mathbf{E}}\) ) and \(x \mapsto i(x)\) is A-linear, it is seen that the external law of composition \((x, u) \mapsto u \llcorner x\) ( \(x \in \mathbf{E}\) , \(u \in \operatorname{Homgr}_{\mathbf{A}}(\mathbf{E}, \mathbf{P})\) ) defines, with addition, a right E-module structure on \(\operatorname{Homgr}_{\mathbf{A}}(\mathbf{E}, \mathbf{P})\) .

In particular we consider the case P = A, \(\operatorname{Homgr}_{\mathbf{A}}(\mathbf{E}, \mathbf{P})\) being in this case the graded dual \(E^{*gr}\) of E; \(i(x)\) is then the graded transpose of the A-linear mapping \(e(x)\) (II, §11, no. 6), in other words, for all x, y in E, \(u \in E^{*gr}\) ,

\[
\langle u \llcorner x, y \rangle = \langle u, x y \rangle .\tag{37}
\]

With the convention at the beginning of the paragraph, note that, if \(x \in E_{p}\) , \(i(x)\) is an endomorphism of \(E^{*gr}\) of degree -p.

For every homogeneous element \(x \in \mathbf{E}_p\) , right multiplication by \(x\) is similarly denoted by \(e'(x)\) and the element \(u \circ e'(x)\) of \(\operatorname{Homgr}_{\mathbf{A}}(\mathbf{E}, \mathbf{P})\) , called the left inner product of \(u\) by \(x\) , by \(x \lrcorner u\) ; we write \(i'(x) = x \lrcorner u\) and \(i'(x)\) is therefore a graded endomorphism of \(\operatorname{Homgr}_{\mathbf{A}}(\mathbf{E}, \mathbf{P})\) of degree \(p\) ; as above this definition can be extended to the case where \(x\) is an arbitrary element of \(\mathbf{E}\) . As in this case \(e'(xy) = e'(y)e'(x)\) ,

\[
i ^ {\prime} (x y) = i ^ {\prime} (x) \circ i ^ {\prime} (y)\tag{38}
\]

which may also be written

\[
x \lrcorner (y \lrcorner u) = (x y) \lrcorner u\tag{39}
\]

and shows that the external law of composition \((x,u)\mapsto x\lrcorner u\) defines, with addition, a left E-module structure on \(\operatorname{Homgr}_{\mathbf{A}}(\operatorname{E},\operatorname{P})\) . The associativity of E implies on the other hand that \(e(x)\circ e'(y)=e'(y)\circ e(x)\) for x,y homogeneous in E, whence the relation

\[
(y \lrcorner u) \llcorner x = y \lrcorner (u \llcorner x)\tag{40}
\]

so that the two external laws of composition on \(\mathrm{Homgr}_{\mathbf{A}}(\mathbf{E},\mathbf{P})\) define on this set an (E, E)-bimodule structure (II, § 1, no. 14).

When we take \(\mathbf{P} = \mathbf{A}\) , \(i'(x)\) is the graded transpose of \(e'(x)\) ; in other words, for all \(x, y\) in \(\mathbf{E}\) , \(u \in \mathbf{E}^{*\mathrm{gr}}\) ,

\[
\langle y, x \lrcorner u \rangle = \langle y x, u \rangle .\tag{41}
\]

When the graded algebra E is commutative, obviously \(u \llcorner x = x \lrcorner u\) . When E is anticommutative and P = A, then for \(x \in E_{p}\) , \(y \in E_{r}\) and \(u \in E_{q}^{*}\) , \(yx = (-1)^{pr}xy\) , whence, by (37) and (41), \(\langle u \llcorner x, y \rangle = (-1)^{pr}\langle y, x \lrcorner u \rangle\) . But as the two sides of this relation are zero except for r = q - p,

\[
x \lrcorner u = (- 1) ^ {p (q - p)} u \llcorner x.
\]

Let F be another graded A-algebra and \(\phi: E \to F\) an A-homomorphism of graded algebras; then \(\tilde{\phi} = \operatorname{Hom}(\phi, 1_{P}): \operatorname{Homgr}_{A}(F, P) \to \operatorname{Homgr}_{A}(E, P)\) is a graded A-homomorphism of degree 0; by definition, for \(x, y\) in E and \(u \in \operatorname{Homgr}_{A}(F, P)\)

\[
\begin{array}{r l} (\tilde {\phi} (u \llcorner \phi (x))) (y) & = (u \llcorner \phi (x)) (\phi (y)) \\ & = u (\phi (x) \phi (y)) = u (\phi (x y)) = (\tilde {\phi} (u)) (x y) = (\tilde {\phi} (u) \llcorner x) (y) \end{array}
\]

or also

\[
\tilde {\phi} (u \llcorner \phi (x)) = \tilde {\phi} (u) \llcorner x\tag{42}
\]

and similarly

\[
\tilde {\phi} (\phi (x) \lrcorner u) = x \lrcorner \tilde {\phi} (u).\tag{43}
\]

In other words, when \(\operatorname{Homgr}_{\mathbf{A}}(\mathbf{F}, \mathbf{P})\) is considered as an (E, E)-bimodule by means of the ring homomorphism \(\phi: \mathbf{E} \to \mathbf{F}\) , it is seen that \(\tilde{\phi}\) is an (E, E)-bimodule homomorphism (or also an E-homomorphism of the (F, F)-bimodule \(\operatorname{Homgr}_{\mathbf{A}}(\mathbf{F}, \mathbf{P})\) into the (E, E)-bimodule \(\operatorname{Homgr}_{\mathbf{A}}(\mathbf{E}, \mathbf{P})\) ).

Examples. In particular the above may be applied when E is one of the graded algebras \(\mathsf{T}(\mathbf{M})\) , \(\mathsf{S}(\mathbf{M})\) or \(\bigwedge(\mathbf{M})\) for an A-module M and P an A-module (with the trivial graduation). To find explicitly the bimodule structures thus obtained, note that the elements of degree -n of \(\operatorname{Homgr}_{\mathbf{A}}(\mathsf{T}(\mathbf{M}),\mathbf{P})\) (resp. \(\operatorname{Homgr}_{\mathbf{A}}(\mathsf{S}(\mathbf{M}),\mathbf{P})\) , resp. \(\operatorname{Homgr}_{\mathbf{A}}(\bigwedge(\mathbf{M}),\mathbf{P})\) ) are identified with the n-linear mappings (resp. symmetric n-linear mappings, resp. alternating n-linear mappings) of \(M^{n}\) into P. It suffices to express the products

\[
f \llcorner (x _ {1} \otimes x _ {2} \otimes \dots \otimes x _ {p})
\]

\((\text{resp. } f \llcorner (x_1 x_2 \ldots x_p), \text{ resp. } f \llcorner (x_1 \wedge x_2 \wedge \dots \wedge x_p))\) for every finite sequence \((x_i)_{1 \leq i \leq p}\) of elements of M and the analogues for the left inner product. It follows immediately from the definitions that

\[
f \llcorner (x _ {1} \otimes x _ {2} \otimes \dots \otimes x _ {p}) = (x _ {1} \otimes x _ {2} \otimes \dots \otimes x _ {p}) \lrcorner f = 0\tag{44}
\]

if \(p > n\) and that, for \(p \leqslant n, f \llcorner (x_1 \otimes x_2 \otimes \dots \otimes x_p)\) (resp.

\[
\left(x _ {1} \otimes x _ {2} \otimes \dots \otimes x _ {p}\right) \lrcorner f)
\]

is the \((n - p)\) -linear mapping defined by

\[
\left\{ \begin{array}{l} (f \llcorner (x _ {1} \otimes x _ {2} \otimes \dots \otimes x _ {p})) (y _ {1}, \ldots , y _ {n - p}) \\ = f (x _ {1}, \ldots , x _ {p}, y _ {1}, \ldots , y _ {n - p}) \\ ((x _ {1} \otimes x _ {2} \otimes \dots \otimes x _ {p}) \lrcorner f) (y _ {1}, \ldots , y _ {n - p}) \\ = f (y _ {1}, \ldots , y _ {n - p}, x _ {1}, \ldots , x _ {p}). \end{array} \right.\tag{45}
\]

For \(p > n\) , there are also in \(\operatorname{Homgr}_{\mathbf{A}}(\mathsf{S}(\mathbf{M}), \mathbf{P})\) (resp. \(\operatorname{Homgr}_{\mathbf{A}}(\bigwedge(\mathbf{M}), \mathbf{P})\) ) formulae (44) with \(x_1 \otimes x_2 \otimes \cdots \otimes x_p\) replaced by \(x_1x_2 \ldots x_p\) (resp. \(x_1 \wedge x_2 \wedge \cdots \wedge x_p\) ). For \(p \leqslant n\) , the same substitutions in (45) define the symmetric \((n - p)\) -linear mappings \(f \llcorner (x_1x_2 \ldots x_p)\) and \((x_1x_2 \ldots x_p) \lrcorner f\) (resp. the alternating \((n - p)\) -linear mappings

\[
f \llcorner (x _ {1} \wedge x _ {2} \wedge \dots \wedge x _ {p}) \quad \text { and } \quad (x _ {1} \wedge x _ {2} \wedge \dots \wedge x _ {p}) \lrcorner f).
\]

When \(n = p\) , the above products are equal to the constant function on \(\mathbf{M}\) equal to \(f(x_{1},\ldots ,x_{p})\) .

If \(u: \mathbf{M} \to \mathbf{N}\) is an A-module homomorphism, \(\mathsf{T}(u): \mathsf{T}(\mathbf{M}) \to \mathsf{T}(\mathbf{N})\) is a graded A-algebra homomorphism, then it follows from what we have seen above that \((\mathsf{T}(u))^{\sim}\) is a \(\mathsf{T}(\mathbf{M})\) -homomorphism of the \((\mathsf{T}(\mathbf{N}), \mathsf{T}(\mathbf{N}))\) -bimodule \(\operatorname{Homgr}_{\mathbf{A}}(\mathsf{T}(\mathbf{N}), \mathbf{P})\) into the \((\mathsf{T}(\mathbf{M}), \mathsf{T}(\mathbf{M}))\) -bimodule \(\operatorname{Homgr}_{\mathbf{A}}(\mathsf{T}(\mathbf{M}), \mathbf{P})\) , relative to the ring homomorphism \(\mathsf{T}(u)\) . There are analogous results for \((\mathbf{S}(u))^{\sim}\) and \((\bigwedge(u))^{\sim}\) .
% semantic-fragments: legacy-input-seam
\relax{}
